\begin{lexlean}{Ledger}
\useglossary{lexlean.std.nat@1.2.0}
\useglossary{spec.titles@1.0.0}
\importmodule{Rows}
\title{Honesty ledger}
\begin{semanticmodule}
definition{
  name: authLeanRel4321,
  parameters: [],
  result: Rows::LedgerClaim,
  body: (record Rows::LedgerClaim {id := "AUTH-LEAN-REL-4-32-1", level := "some-true", statement := "Lean 4.32.1 is the release tagged leanprover/lean4:v4.32.1, built from source commit f054605aea4b840552cca2e725580bffd1e1b704.", authority := "LEAN-REL-4-32-1"}),
  executable: true
}

definition{
  name: authLake4321,
  parameters: [],
  result: Rows::LedgerClaim,
  body: (record Rows::LedgerClaim {id := "AUTH-LAKE-4-32-1", level := "some-true", statement := "Lake as distributed with Lean 4.32.1 resolves a pinned workspace environment through lake env without network access when every locked dependency is locally available.", authority := "LAKE-4-32-1"}),
  executable: true
}

definition{
  name: authLeanchecker4321,
  parameters: [],
  result: Rows::LedgerClaim,
  body: (record Rows::LedgerClaim {id := "AUTH-LEANCHECKER-4-32-1", level := "some-true", statement := "The leanchecker executable distributed with Lean 4.32.1 replays a compiled module's environment through the Lean kernel in a separate process and is not an independent proof checker.", authority := "LEANCHECKER-4-32-1"}),
  executable: true
}

definition{
  name: authLeanLcnf4321,
  parameters: [],
  result: Rows::LedgerClaim,
  body: (record Rows::LedgerClaim {id := "AUTH-LEAN-LCNF-4-32-1", level := "some-true", statement := "Lean 4.32.1 translates a code-generating definition to base-phase LCNF with Lean.Compiler.LCNF.toDecl, the compiler's own input to its later passes, so every constant the compiled definition can depend on at run time is named by that translation or by a definition it names.", authority := "LEAN-LCNF-4-32-1"}),
  executable: true
}

definition{
  name: authRustc1971,
  parameters: [],
  result: Rows::LedgerClaim,
  body: (record Rows::LedgerClaim {id := "AUTH-RUSTC-1-97-1", level := "some-true", statement := "rustc 1.97.1 compiles a safe Rust 2021 crate with the semantics The Rust Reference states for that release: the checked integer methods return None exactly on overflow, the wrapping shifts mask their amount, integer division truncates toward zero, and a #![no_std] crate that declares no extern crate links neither alloc nor std.", authority := "RUSTC-1-97-1"}),
  executable: true
}

definition{
  name: authPrintAxioms4321,
  parameters: [],
  result: Rows::LedgerClaim,
  body: (record Rows::LedgerClaim {id := "AUTH-PRINT-AXIOMS-4-32-1", level := "some-true", statement := "Lean 4.32.1 reports the transitive axiom dependencies of a declaration through #print axioms in exactly the output forms recorded by the committed axiom-parser fixtures.", authority := "PRINT-AXIOMS-4-32-1"}),
  executable: true
}

definition{
  name: authUorGnaf1Draft2,
  parameters: [],
  result: Rows::LedgerClaim,
  body: (record Rows::LedgerClaim {id := "AUTH-UOR-GNAF-1-DRAFT-2", level := "some-true", statement := "UOR-GNAF normative draft uor-gnaf/1-draft.2 requires an optimality claim's machine contract, accounting model, complete-system universe and its optimizer-independent completeness evidence, objective order, and claim class to be fixed before optimization, and treats UOR-NAF as informative only.", authority := "UOR-GNAF-1-DRAFT-2"}),
  executable: true
}

definition{
  name: authNistFips1804,
  parameters: [],
  result: Rows::LedgerClaim,
  body: (record Rows::LedgerClaim {id := "AUTH-NIST-FIPS-180-4", level := "some-true", statement := "SHA-256 is the function FIPS 180-4 specifies in its section 6.2, defined for messages of fewer than 2^64 bits.", authority := "NIST-FIPS-180-4"}),
  executable: true
}

definition{
  name: authNistSha256Examples,
  parameters: [],
  result: Rows::LedgerClaim,
  body: (record Rows::LedgerClaim {id := "AUTH-NIST-SHA256-EXAMPLES", level := "some-true", statement := "NIST's SHA-256 examples give ba7816bf8f01cfea414140de5dae2223b00361a396177a9cb410ff61f20015ad as the digest of abc and 248d6a61d20638b8e5c026930c3e6039a33ce45964ff2167f6ecedd419db06c1 as the digest of the 448-bit message abcdbcdecdefdefgefghfghighijhijkijkljklmklmnlmnomnopnopq.", authority := "NIST-SHA256-EXAMPLES"}),
  executable: true
}

definition{
  name: authNistFips1802,
  parameters: [],
  result: Rows::LedgerClaim,
  body: (record Rows::LedgerClaim {id := "AUTH-NIST-FIPS-180-2", level := "some-true", statement := "FIPS 180-2 gives cdc76e5c9914fb9281a1c7e284d73e67f1809a48a497200e046d39ccc7112cd0 as the SHA-256 digest of one million repetitions of a.", authority := "NIST-FIPS-180-2"}),
  executable: true
}

definition{
  name: authNistCavpSha256Shortmsg,
  parameters: [],
  result: Rows::LedgerClaim,
  body: (record Rows::LedgerClaim {id := "AUTH-NIST-CAVP-SHA256-SHORTMSG", level := "some-true", statement := "The NIST CAVP file SHA256ShortMsg.rsp states 65 SHA-256 message and digest pairs, of lengths 0 to 512 bits in multiples of 8, among them e3b0c44298fc1c149afbf4c8996fb92427ae41e4649b934ca495991b7852b855 as the digest of the empty message.", authority := "NIST-CAVP-SHA256-SHORTMSG"}),
  executable: true
}

definition{
  name: authNistCavpSha256Longmsg,
  parameters: [],
  result: Rows::LedgerClaim,
  body: (record Rows::LedgerClaim {id := "AUTH-NIST-CAVP-SHA256-LONGMSG", level := "some-true", statement := "The NIST CAVP file SHA256LongMsg.rsp states 64 SHA-256 message and digest pairs, of lengths 1304 to 51200 bits.", authority := "NIST-CAVP-SHA256-LONGMSG"}),
  executable: true
}
\end{semanticmodule}
\end{lexlean}
